% Edoardo Galli, CV. One A4 page, built with LuaLaTeX:
%   latexmk -lualatex -outdir=<build dir> <this file>
%
% Layout: a 3.6cm rail on the left holds dates, section headings and labels;
% a 13.4cm content column holds everything else. The text block is the
% content column, and \rail hangs its argument into the left margin, so the
% PDF text keeps rail, title and bullets in reading order for each entry.
%
% Type system: Inter Regular and SemiBold only, three sizes (22pt name,
% 10pt body, 9pt meta), three colours (ink, muted grey, one accent bar).
% Every bullet is set in an unbreakable box: a bullet too long for its line
% raises an overfull box warning instead of wrapping.

\documentclass[a4paper]{article}

% 2cm page margin + 3.6cm rail on the left; 2cm on the right.
\usepackage[a4paper, left=5.6cm, right=2cm, top=2.2cm, bottom=1.8cm]{geometry}
\usepackage{fontspec}
\usepackage{xcolor}
\usepackage[hidelinks]{hyperref}

\hypersetup{
  pdftitle  = {Edoardo Galli, CV},
  pdfauthor = {Edoardo Galli},
}

% ---------------------------------------------------------------- Colour --

\definecolor{ink}{HTML}{171717}    % all text and links
\definecolor{muted}{HTML}{606060}  % dates, meta lines, bullets
\definecolor{accent}{HTML}{E8A33D} % the bar above the name; never text

% ------------------------------------------------------------------ Type --

% LetterSpace is a percentage of the em: 1.0 = +10/1000 em.
% Tabular figures stay in the rail: Inter's tnum also widens the hyphen.
\defaultfontfeatures{
  Extension    = .otf,
  UprightFont  = *-Regular,
  BoldFont     = *-SemiBold,
  Numbers      = Lining,
}
\setmainfont{Inter}[LetterSpace = 1.0]                % body, +10/1000 em
\newfontfamily\metafamily{Inter}[LetterSpace = 1.6]   % meta, +16/1000 em
\newfontfamily\namefamily{Inter}[LetterSpace = -1.8]  % name, -18/1000 em

\newcommand\namestyle{\namefamily\bfseries\fontsize{22pt}{24.2pt}\selectfont}
\newcommand\bodystyle{\normalfont\fontsize{10pt}{13.5pt}\selectfont}
\newcommand\metastyle{\metafamily\fontsize{9pt}{12pt}\selectfont}

% ---------------------------------------------------------------- Layout --

\newlength\railwidth   \setlength\railwidth{3.6cm}
\newlength\sectiongap  \setlength\sectiongap{14pt}
\newlength\headinggap  \setlength\headinggap{2pt}
\newlength\entrygap    \setlength\entrygap{6pt}
\newlength\bulletgap   \setlength\bulletgap{1pt}
\newlength\bulletindent\setlength\bulletindent{10pt}
\newlength\contactgap  \setlength\contactgap{1.2em}

\pagestyle{empty}
\setlength\parindent{0pt}
\setlength\parskip{0pt}
\frenchspacing
\raggedright

% \rail{<text>}: grey meta text set in the rail, left of the current line.
\newcommand\rail[1]{%
  \leavevmode\llap{\makebox[\railwidth][l]{%
    \metastyle\addfontfeatures{Numbers = Tabular}\color{muted}#1}}}

% \range{<from>}{<to>}: a date range, en dash between thin spaces.
\newcommand\range[2]{#1\,–\,#2}

% railed: paragraphs that start at the rail edge and may use the full 17cm,
% for the header and the section headings.
\newenvironment{railed}{\par\leftskip=-\railwidth}{\par}

% \cvsection{<heading>}: a SemiBold heading in the rail, on its own line.
\newif\ifsectionstart
\newcommand\cvsection[1]{%
  \par\vspace{\sectiongap}%
  \begin{railed}\bfseries #1\end{railed}%
  \nobreak\vspace{\headinggap}\sectionstarttrue}

% \entry{<rail: dates>}{<title line>}; entries after the first in a section
% sit \entrygap below the one before.
\newcommand\entry[2]{%
  \par\ifsectionstart\sectionstartfalse\else\addvspace{\entrygap}\fi
  \rail{#1}#2\par}

% \meta{<text>}: a grey 9pt line in the content column.
\newcommand\meta[1]{{\metastyle\color{muted}#1\par}}

% \point{<text>}: a one-line bullet; the grey dot sits on the content edge.
\newcommand\point[1]{%
  \par\addvspace{\bulletgap}%
  \makebox[\bulletindent][l]{\color{muted}•}\mbox{#1}\par}

% \skill{<label>}{<items>}: label in the rail, items in the content column.
\newcommand\skill[2]{\par\addvspace{\bulletgap}\rail{#1}#2\par}

\begin{document}
\bodystyle\color{ink}

% ---------------------------------------------------------------- Header --

\begin{railed}
  {\color{accent}\rule{1.2cm}{1.5pt}}\par\vspace{10pt}
  {\namestyle Edoardo Galli\par}\vspace{3pt}
  {\color{muted}Software engineer building AI agents, based in Bologna\par}
  \vspace{5pt}
  {\metastyle
    \href{mailto:edoski.dev@gmail.com}{edoski.dev@gmail.com}\hspace{\contactgap}%
    \href{https://edoski.com}{edoski.com}\hspace{\contactgap}%
    \href{https://github.com/edoski}{github.com/edoski}\hspace{\contactgap}%
    \href{https://www.linkedin.com/in/edoardo-galli-5074321b9}{linkedin.com/in/edoardo-galli-5074321b9}\par}
\end{railed}

% ------------------------------------------------------------ Experience --

\cvsection{Experience}

\entry{\range{Sep 2026}{Present}}{\textbf{Software Engineer}, Amoreg}
\point{Lead development of Iris, an AI agent for safety-gated LinkedIn outreach}
\point{Ship features on Amoreg's core product, an agentic compliance platform}
\point{Build internal tooling, agent infrastructure and CI/CD for the team}

\entry{\range{Apr}{Sep 2026}}{\textbf{Research Intern}, DISI, University of Bologna}
\meta{Supervised by Prof. Federico Montori and Dr. Ivan Zyrianoff}
\point{Researched deep learning for EVM transaction timing}
\point{Built \href{https://github.com/edoski/kairos}{KAIROS} end to end: data pipelines, HPC training, on-device inference}

% ------------------------------------------------------------- Education --

\cvsection{Education}

\entry{\range{Sep 2026}{Present}}{\textbf{M.Sc. Artificial Intelligence}, University of Bologna}

\entry{\range{Sep 2021}{Oct 2026}}{\textbf{B.Sc. Information Science for Management}, University of Bologna}
\meta{Final grade 108/110}

% ----------------------------------------------------------- Publication --

\cvsection{Publication}

% The title is too long for one line; break it at the phrase boundary.
\entry{2027}{\textbf{KAIROS: A Predictive Framework for Cost Optimization}\\
  \textbf{in Blockchain Environments}}
I. Zyrianoff, \textbf{E. Galli}, A. Esposito, L. Gigli, M. Di Felice and F. Montori\par
\meta{IEEE Consumer Communications \& Networking Conference (CCNC 2027)}

% -------------------------------------------------------------- Projects --

\cvsection{Projects}

\entry{}{\href{https://github.com/edoski/servatus}{\textbf{Servatus}}\hspace{\contactgap}Zero-dependency PyPI package for resumable Slurm campaigns}
\meta{Python, Slurm, OpenSSH, Apptainer, POSIX}

\entry{}{\href{https://github.com/edoski/journal}{\textbf{Journal}}\hspace{\contactgap}Automates Flow focus data, goals and reminders in Obsidian}
\meta{Python, Pydantic, pytest, Obsidian}

% ---------------------------------------------------------------- Skills --

\cvsection{Skills}

\skill{AI agents}{Claude Agent SDK, MCP, Vercel AI SDK, model evaluation}
\skill{Machine learning}{Deep learning, time series forecasting, PyTorch Lightning, ExecuTorch}
\skill{Data and HPC}{Python, Polars, pandas, Slurm, Apptainer, Docker, Linux}
\skill{Web}{TypeScript, React, Next.js, PostgreSQL, Tailwind CSS}
\skill{Languages}{English (IELTS Academic 8.5), Italian (native)}

\end{document}
